\subsection{Conformal scaling identity}

If $g=\rho h$ on a surface, then $g^{-1}=\rho^{-1}h^{-1}$.  Therefore, for any
smooth real function $a$,
\begin{equation}
 |\nabla a|_g^2=|\dd a|_{g^{-1}}^2
 =\rho^{-1}|\dd a|_h^2.
 \label{eq:appendix-conformal-scaling}
\end{equation}
Taking
$\rho=|\dd a|_h^2/(\mu^2+\lambda^2a^2)$ proves
Theorem~\ref{thm:conformal-realization}.  The construction is valid exactly where
$\dd a\ne0$; at a critical point the displayed coefficient vanishes and is not a
Riemannian metric.

For $a_h=e^xh(y)$,
\[
 \partial_xa_h=e^xh,
 \qquad
 \partial_ya_h=e^xh',
\]
which gives Equation~\eqref{eq:parallel-metric}.  For the helical assignment,
writing $\xi=ny-m\theta$ gives
\[
 \partial_xA=ne^{nx}\sin\xi,
 \qquad
 \partial_yA=ne^{nx}\cos\xi,
\]
so $|\dd^VA|^2=n^2e^{2nx}$ independently of $\xi$.  This cancellation is the
reason the helical metric has the particularly simple numerator in
\eqref{eq:helical-metric}.

\subsection{Vertical transversality in coordinates}

Locally trivialize $q:X\to B$ and $E\to X$ near a zero so that
$X\cong U\times V$ and $E\cong(U\times V)\times\R^r$, with $U\subset\R^m$ a
fiber chart and $V\subset\R^k$ a base chart.  The section becomes a map
$s(u,b)\in\R^r$.  Vertical transversality says that the $r\times m$ matrix
$D_us$ has rank $r$.  The full Jacobian $(D_us,D_bs)$ then has rank $r$, proving
that $s^{-1}(0)$ has dimension $m+k-r$.

To see the projection submersion explicitly, solve
\[
 D_us\,\dot u=-D_bs\,\dot b
\]
for $\dot u$ using surjectivity of $D_us$.  Then $(\dot u,\dot b)$ is tangent to
the incidence and projects to the arbitrary vector $\dot b$.  This is the local
linear-algebra content of Theorem~\ref{thm:parameterized-zero-section}.

\subsection{Properness checks}

For the helical model, $M_L\times S^1$ is compact and the zero incidence is the
inverse image of the closed set $\{0\}$ under a continuous function.  It is
therefore compact, and every map from it to the Hausdorff circle is proper.  At
$x=\pm L$, the vertical normal to the boundary is $\partial_x$, while
$\partial_yA\ne0$ on the zero set.  Hence the zero incidence is neat and its
boundary projects submersively.

For the escape family of Proposition~\ref{prop:nonproper-escape}, fix any
$\varepsilon>0$ and set $t_j=\varepsilon/(j+1)$.  The sequence
\[
 (x_j,y_j,t_j)=(0,t_j^{-1},t_j)
\]
lies in the zero incidence above $[0,\varepsilon]$ and leaves every compact subset
of $\R^2\times[0,\varepsilon]$.  This is a direct certificate that the incidence
projection is not proper.

\subsection{Audit records for the explicit models}

The following records make all data and limitations explicit.

\paragraph{Cylindrical product model.}
Domain $\R\times S^1$; metric $\dd r^2+\dd\phi^2$; assignment
$(\mu/\lambda)\sinh(\lambda(r-c))$ for $\lambda\ne0$ and $\mu(r-c)$ for
$\lambda=0$; singular set empty; equation verified by differentiation; zero
topology one circle; compact parameter base gives a proper product incidence;
status: proved.

\paragraph{Parallel multi-zero model.}
Domain $\R^2$; metric \eqref{eq:parallel-metric}; assignment $e^xh(y)$; singular
set empty when all zeros of $h$ are simple; equation verified by conformal
realization; zero topology one line for each zero of $h$; no parameter required;
completeness not asserted; status: proved, newly constructed in this paper.

\paragraph{Logarithmic model.}
Downstairs domain $\C^\times$, upstairs domain $\C$; metrics and assignments in
\eqref{eq:punctured-plane-aes} and \eqref{eq:logarithmic-aes}; singular set empty;
equation verified by pullback; two zero components downstairs and countably many
upstairs; deck action $k\mapsto k+2$; missing origin removable; status: proved
covering model, not an essential singular AES.

\paragraph{Holomorphic branched pullback.}
Domain a Riemann surface $\Sigma$; map a nonconstant holomorphic
$\Phi:\Sigma\to\C$; assignment $\im\Phi$; metric
\eqref{eq:holomorphic-pullback-metric}; singular set $\Crit(\Phi)$; equation
verified by pullback from the linear target AES; a degree-$m$ critical zero has
$2m$ prongs and a $2\pi m$ metric cone; status: proved local theorem, with no
claim that an arbitrary $\Phi$ is history-natural.

\paragraph{Finite Blaschke bridge.}
Domain $\mathbb D$; map a finite Blaschke product $B$; finite divisor
$B^{-1}(0)$; phase assignments and metrics as in
Example~\ref{ex:blaschke-bridge}; singular set $\Crit(B)$ for the phase model and
$B^{-1}(0)\cup\Crit(B)$ for the pulled-back radial disc model; zero topology a
finite holomorphic divisor or the preimage of a real diameter, depending on field
rank; status: proved finite bridge, but its freely selected roots are not a
history-derived invariant.

\paragraph{Essential sine control.}
Domain a disc with declared singular set
$\{0\}\cup\{(\pi/2+k\pi)^{-1}\}$; assignment
$\im\sin(1/z)$ and metric in Example~\ref{ex:sin-essential-zero-network}; zero
topology one real-axis branch and countably many circles, with four-pronged
$4\pi$ cone nodes accumulating at the essential singularity; status: proved
analytic-capacity example, not the arithmetic model.

\paragraph{$q=4$ Hecke--Hauptmodul model.}
Domain $\HH$; normalized Hauptmodul $\beta$ with orbifold signature
$(2,4,\infty)$, finite ramification indices $2$ and $4$, and a simple pole at the
cusp; singular set $\mathcal E_4=\beta^{-1}(1)$; assignment and metric
\eqref{eq:q4-aeg-data}; equation verified in the local coordinate $\log W$;
zero topology $\beta^{-1}([0,1])$, which becomes the $\{\infty,4\}$ Hecke graph
after bivalent vertices are suppressed; order-$4$ nodes have four prongs and cone
angle $4\pi$ upstairs; the metric descends while the assignment descends as a
sign-line-bundle section; regular locus incomplete at the cone nodes, completed
full quotient a complete length space with its cusp at infinite distance; status:
proved from the normalized automorphic quotient, at the projective-operator history
level only.

\paragraph{Helical model.}
Domain $[-L,L]\times S^1$; parameter $\theta\in S^1$; metric
\eqref{eq:helical-metric}; assignment \eqref{eq:helical-assignment}; singular set
empty; equation verified above; each zero fiber has $2n$ intervals; total incidence
has $2\gcd(n,m)$ annuli and is proper with neat boundary; status: proved.

\paragraph{Morse models.}
Domain $\R^2$, singular set $\{0\}$, parameter $\tau\in\R$; metric
\eqref{eq:morse-aeg-metric}; assignments \eqref{eq:morse-assignments}; equation
verified on the punctured plane; definite birth/death and indefinite reconnection;
the total zero incidence is smooth but its parameter projection is critical;
status: proved examples, not universal AEG normal forms.

\paragraph{Complex branch model.}
Background is the basic hyperbolic AES from Paper II with $\mu,\lambda>0$; field
$F_\tau(w)=w^2-\tau$ is complex-valued and arithmetic holomorphic in $w$; the
field is not the real assignment; root incidence is smooth, projection has one
simple branch point; boundary gives a half-twist; status: proved local root model.

\subsection{Completeness warning}

Conformal realization does not preserve completeness.  If a conformal coefficient
$\rho$ decays rapidly along an end, an $h$-infinite path may have finite
$g=\rho h$ length.  In the parallel examples, the factor $e^{2x}$ commonly makes
$x\to-\infty$ a finite-distance end.  Any theorem that needs complete fibers must
therefore add completeness as a hypothesis and recheck each model; it cannot infer
it from the eikonal identity.

The Hecke model illustrates the required distinction.  Its regular locus is
incomplete because the deleted order-$4$ points lie at finite distance.  The
singular length completion adjoins $4\pi$ cone points, while the cusp coordinate
has metric asymptotic to $|\dd q/q|^2$ and is at infinite distance.  This audited
completion does not turn a cone point into a smooth regular AES point.
